% ============================================================
\section{Arquitectura del \emph{pipeline}}
\label{sec:arquitectura}
% ============================================================

El \emph{pipeline} vive en el repositorio \archivo{smn_toe}, tiene un único punto de entrada (\texttt{run.sh}) y deja sus salidas en \texttt{data/processed/}, de donde el Entregable~2 las lee sin necesidad de conocer el detalle interno de este módulo. Se rige por tres principios:

\begin{itemize}
\item \textbf{Idempotencia.} Cada paso omite el procesamiento cuya salida ya existe, de modo que una corrida puede interrumpirse y reanudarse sin recalcular lo hecho. La excepción deliberada es el control de calidad, que se recalcula completo en cada corrida porque es barato y un resultado en caché puede quedar obsoleto.
\item \textbf{Verificación por valor y no por metadato.} Las unidades, los valores de relleno y la lista de meses se comprueban sobre los datos numéricos y no sobre los atributos declarados en los archivos. Existen archivos cuyo atributo contradice su contenido; por ejemplo, archivos que declaran grados Celsius y contienen valores en Kelvin.
\item \textbf{Cascada de fuentes con degradación explícita.} Si el nodo principal de ESGF no resuelve un modelo, se prueban automáticamente nodos alternativos y Copernicus CDS. Lo que ninguna fuente resuelve queda registrado en un reporte de estado, en lugar de fallar en silencio.
\end{itemize}

\texttt{run.sh} no contiene lógica de procesamiento propia: ejecuta los pasos 00, 00b, 01, 02, 03, 04, 05, 06 y 07 en ese orden fijo, o solo un rango de ellos, y registra la salida de cada uno en \texttt{logs/pipeline.log}. Al terminar cualquier corrida, completa o parcial, genera siempre el registro de modelos, las figuras automáticas, el reporte de estado y un paquete comprimido con los artefactos para actualizar el informe. Así, cada ejecución deja constancia verificable de qué modelos quedaron descargados, procesados o pendientes. La instalación, los argumentos de \texttt{run.sh}, los tiempos y el espacio en disco requeridos, y la interpretación de las salidas se describen en el documento de anexos que acompaña a este informe.
